\begin{abstract}
\noindent
Fungal biosynthetic gene clusters (BGCs) underlie a chemically and biologically important class of
secondary metabolites, yet the primary evidence for any given cluster --- which genes were actually
shown to be required, which metabolites were directly observed, under what conditions, and how
confidently --- is scattered across decades of literature and only partially captured by curated
sequence repositories such as MIBiG. We present the \textbf{Fungal BGC Atlas}, a curated,
evidence-linked knowledge base of 609 fungal BGC dossiers, each anchored to a MIBiG accession and
built by large-language-model-assisted extraction from primary literature followed by full manual
review of every dossier. Beyond the cluster-level record, the Atlas resolves each dossier into 1572
genes, 807 metabolites, 915 experiments, 936 measurements, 1534 claims and 1264 curated
subject--predicate--object relationships, all traceable to 2575 source records through 47\,768
explicit evidence links. Structural and functional context is captured along 27 independent facet
axes (biosynthesis class, chemistry, ecology, genome context and activation conditions), and 995
conflicts or unresolved uncertainties are preserved as first-class records rather than collapsed
into a forced consensus. We package the Atlas as a relational SQLite database, a knowledge graph,
and a self-contained interactive web dashboard (Overview, Taxonomy, Chemistry, Facet Explorer,
Knowledge Graph, Evidence, Conflicts, a searchable Dossier Browser with cross-table detail views, and
a live Data Dictionary), deployed on GitHub Pages and archived on Zenodo with an open licence
(CC-BY-4.0). The Atlas turns 609 individually FAIR-adjacent literature records into a single
queryable, comparable, and citable resource for fungal natural-product genome mining.\\[0.5em]
\textbf{Database URL:} \url{https://dr-richard-barker.github.io/fungal-bgc-atlas/}
\end{abstract}

\textbf{Keywords:} biosynthetic gene clusters; fungal natural products; MIBiG; knowledge graph;
FAIR data; curated database; secondary metabolism; evidence provenance
